\chapter{Appendix C: Installation Guide and Operations Runbook}

\section{Scope}
This appendix provides a complete practical guide to install, configure, run, troubleshoot, and maintain the SmartHome platform in development and pilot production environments.

\section{Prerequisites}
\subsection{Software}
\begin{itemize}[leftmargin=1.2cm]
  \item Node.js LTS (recommended v20+).
  \item npm (comes with Node.js).
  \item MongoDB server (local or managed).
  \item Optional: MQTT client tools for diagnostics.
  \item Optional: PM2 or container runtime for service supervision.
\end{itemize}

\subsection{System}
\begin{itemize}[leftmargin=1.2cm]
  \item Minimum 4 GB RAM for development.
  \item Open ports for backend HTTP and MQTT.
  \item Stable network for OAuth/email features.
\end{itemize}

\section{Repository Setup}
\begin{lstlisting}[style=codeblock,basicstyle=\ttfamily\small]
git clone <repository-url>
cd SmartHome-PFE
npm install
cd backend && npm install
cd ../frontend && npm install
\end{lstlisting}

\section{Environment Configuration}
Create backend environment file with keys such as:
\begin{lstlisting}[style=codeblock]
PORT=5000
MQTT_PORT=1883
MONGODB_URI=mongodb://127.0.0.1:27017/SmartHome
JWT_SECRET=replace_with_secure_secret
FRONTEND_URL=http://localhost:5173
BACKEND_URL=http://localhost:5000
ADMIN_ACCESS_CODE=ADMIN123
AGENCY_ACCESS_CODE=AGENCY2024
GOOGLE_CLIENT_ID=...
GOOGLE_CLIENT_SECRET=...
\end{lstlisting}

Frontend environment sample:
\begin{lstlisting}[style=codeblock]
VITE_API_BASE_URL=http://localhost:5000/api
\end{lstlisting}

\section{Local Run Commands}
\subsection{Monorepo Orchestrated Run}
\begin{lstlisting}[style=codeblock]
npm run start
\end{lstlisting}

\subsection{Separate Run}
\begin{lstlisting}[style=codeblock]
# terminal 1
cd backend
npm run dev

# terminal 2
cd frontend
npm run dev
\end{lstlisting}

\section{Health Verification Checklist}
\begin{enumerate}[leftmargin=1.2cm]
  \item Backend logs show successful MongoDB connection.
  \item HTTP API responds on configured port.
  \item Frontend starts and route guard logic works.
  \item MQTT TCP port accepts connections.
  \item Socket authenticated events delivered after login.
\end{enumerate}

\section{First-use Functional Smoke Tests}
\subsection{Auth Smoke}
\begin{itemize}[leftmargin=1.2cm]
  \item Register resident and login.
  \item Register or login admin.
  \item Confirm token-based API access.
\end{itemize}

\subsection{Device Smoke}
\begin{itemize}[leftmargin=1.2cm]
  \item Create a device as admin.
  \item Send a command to the device endpoint.
  \item Confirm MQTT publish and DB state update.
\end{itemize}

\subsection{Permission Smoke}
\begin{itemize}[leftmargin=1.2cm]
  \item Submit resident permission request.
  \item Review and approve from admin session.
  \item Confirm resident receives realtime status update.
\end{itemize}

\subsection{Gas Smoke}
\begin{itemize}[leftmargin=1.2cm]
  \item Publish gas reading below threshold.
  \item Publish reading above threshold.
  \item Confirm emergency flag and emitted alert.
\end{itemize}

\section{MQTT Diagnostic Recipes}
\subsection{Subscribe to command channel}
\begin{lstlisting}[style=codeblock]
mosquitto_sub -h localhost -p 1883 -t "command/#" -v
\end{lstlisting}

\subsection{Publish gas reading}
\begin{lstlisting}[style=codeblock]
mosquitto_pub -h localhost -p 1883 -t "sensors/H1234/gas" -m "{\"level\":650}"
\end{lstlisting}

\section{Common Incidents and Resolution}
\subsection{Incident: Cannot find module passport}
\textbf{Symptoms}: backend crashes at startup with module-not-found.

\textbf{Resolution}:
\begin{enumerate}[leftmargin=1.2cm]
  \item Verify backend dependencies installed.
  \item Install missing package in backend project.
  \item Restart backend runtime.
\end{enumerate}

\subsection{Incident: Socket authentication failed}
\textbf{Symptoms}: client cannot receive realtime updates.

\textbf{Resolution}:
\begin{enumerate}[leftmargin=1.2cm]
  \item Ensure valid JWT is sent in socket auth payload.
  \item Verify JWT secret consistency across environments.
  \item Confirm user still exists and token is not expired.
\end{enumerate}

\subsection{Incident: MQTT commands not reaching devices}
\textbf{Resolution actions}:
\begin{itemize}[leftmargin=1.2cm]
  \item Validate broker port and topic naming.
  \item Check payload serialization and parser logic.
  \item Confirm device simulator/client subscribes expected pattern.
\end{itemize}

\section{Security Hardening Checklist}
\begin{itemize}[leftmargin=1.2cm]
  \item Replace all default secrets and access codes.
  \item Enforce HTTPS and secure proxy headers.
  \item Restrict CORS origins.
  \item Add request rate limiting.
  \item Add audit logs for admin actions.
  \item Introduce secret rotation process.
\end{itemize}

\section{Backup and Restore Procedure}
\subsection{Backup}
\begin{lstlisting}[style=codeblock]
# example
mongodump --uri="mongodb://127.0.0.1:27017/SmartHome" --out=./backup
\end{lstlisting}

\subsection{Restore}
\begin{lstlisting}[style=codeblock]
# example
mongorestore --uri="mongodb://127.0.0.1:27017/SmartHome" ./backup/SmartHome
\end{lstlisting}

\section{Release Management}
\begin{enumerate}[leftmargin=1.2cm]
  \item Create release branch and freeze scope.
  \item Execute test suite and smoke checklist.
  \item Tag release candidate and deploy to staging.
  \item Validate monitoring and rollback readiness.
  \item Promote to production with changelog.
\end{enumerate}

\section{Rollback Strategy}
\begin{itemize}[leftmargin=1.2cm]
  \item Keep prior deploy artifacts accessible.
  \item Version database migrations and backups.
  \item Revert frontend static build and backend image atomically.
\end{itemize}

\section{Service Ownership Matrix}
\begin{table}[H]
\centering
\caption{Operational ownership examples}
\begin{tabular}{p{5cm}p{4cm}p{5cm}}
\toprule
Component & Primary Owner & Secondary Owner \\
\midrule
Backend API & Backend Engineer & Full-stack Engineer \\
Frontend SPA & Frontend Engineer & Full-stack Engineer \\
MongoDB & DevOps Engineer & Backend Engineer \\
MQTT Broker & IoT Engineer & Backend Engineer \\
Incident Response & DevOps Lead & Product Tech Lead \\
\bottomrule
\end{tabular}
\end{table}

\section{Capacity Planning Notes}
For scaling households and devices, monitor:
\begin{itemize}[leftmargin=1.2cm]
  \item peak concurrent socket connections,
  \item MQTT message rate,
  \item DB write amplification during burst events,
  \item API auth traffic under login spikes.
\end{itemize}

\section{Production Readiness Checklist}
\begin{itemize}[leftmargin=1.2cm]
  \item Secrets externalized and rotated.
  \item Monitoring and alerting active.
  \item Backup/restore tested.
  \item Security scan completed.
  \item Runbooks approved.
\end{itemize}

\section{Appendix Summary}
This runbook transforms the project from academic implementation to maintainable operational service by documenting reproducible procedures, controls, and recovery practices.
